\begin{titlepage}
\begin{tikzpicture}[remember picture,overlay]
  \fill[walmartnavy] (current page.north west) rectangle ([yshift=-4.2cm]current page.north east);
  \fill[walmartyellow] ([yshift=-4.2cm]current page.north west) rectangle ([yshift=-4.35cm]current page.north east);
  \fill[softblue] ([yshift=1.15cm]current page.south west) rectangle (current page.south east);
  \foreach \a in {0,60,...,300}{
    \fill[walmartyellow,rotate around={\a:(15.5,25.2)}] (15.5,25.2) ellipse (0.15 and 0.55);
  }
\end{tikzpicture}

\vspace*{0.3cm}
\begin{center}
{\color{white}\Large\bfseries \placeholder{nom de l'établissement}}\\[0.2cm]
{\color{white}\normalsize \placeholder{filière et spécialité}}\\[2.4cm]

{\Large\textsc{Rapport de projet de fin d'études}}\\[0.9cm]
{\Huge\bfseries\color{walmartnavy} Walmart Streaming Flow}\\[0.35cm]
{\Large\bfseries Conception et mise en oeuvre d'une plateforme Data Engineering incrémentale pour l'analyse retail}\\[0.8cm]

\fcolorbox{walmartblue}{softblue}{%
\parbox{0.84\textwidth}{\centering
\large PostgreSQL \(\rightarrow\) Databricks Lakehouse \(\rightarrow\) dbt \(\rightarrow\) Power BI\\
\normalsize Génération continue d'événements, traitement horaire et modèle décisionnel en étoile}}

\vfill
\begin{tabularx}{0.9\textwidth}{>{\bfseries}l X}
Réalisé par : & Amine Ait Ali \\
Type de projet : & \placeholder{PFE, stage ou projet académique à confirmer} \\
Organisme d'accueil : & \placeholder{nom de l'entreprise ou du laboratoire} \\
Encadrant académique : & \placeholder{nom et fonction} \\
Encadrant professionnel : & \placeholder{nom et fonction, si applicable} \\
Période : & \placeholder{dates de début et de fin} \\
Année universitaire : & \placeholder{année universitaire} \\
\end{tabularx}
\vspace{0.6cm}
{\small\textit{Projet pédagogique et de portfolio, sans affiliation à Walmart Inc.}}
\end{center}
\end{titlepage}

\chapter*{Remerciements}
\addcontentsline{toc}{chapter}{Remerciements}
Nous tenons à exprimer notre gratitude à \placeholder{nom de l'encadrant académique} pour son accompagnement, ses orientations méthodologiques et la qualité de ses retours durant la réalisation de ce projet. Nous remercions également \placeholder{nom de l'encadrant professionnel ou de la structure d'accueil} pour \placeholder{nature de l'accompagnement et ressources mises à disposition}.

Nous adressons nos remerciements aux enseignants de \placeholder{nom de la formation} qui nous ont transmis les fondements nécessaires en bases de données, ingénierie des données, systèmes distribués et informatique décisionnelle. Enfin, nous remercions nos proches pour leur soutien constant. Les choix techniques, analyses et éventuelles erreurs présentés dans ce document restent de notre responsabilité.

\clearpage
\chapter*{Résumé}
\addcontentsline{toc}{chapter}{Résumé}
Ce rapport présente la conception et la réalisation de \emph{Walmart Streaming Flow}, une plateforme pédagogique de Data Engineering destinée à simuler, intégrer, transformer et valoriser des événements retail. Un générateur web en Python et Flask produit deux familles d'événements cohérents avec le modèle métier : l'inscription automatique de nouveaux clients et la création de commandes comportant une à cinq lignes. Les écritures sont réalisées de manière transactionnelle dans une base PostgreSQL externe, tout en préservant les clés étrangères reliant clients, magasins, employés et produits.

Le traitement analytique demeure volontairement micro-batch et horaire. Apache Airflow orchestre onze étapes déterministes : ingestion incrémentale vers Bronze, contrôles de fraîcheur, transformations dbt vers Silver technique puis Silver métier, snapshots historisés et construction de la couche Gold. Le lakehouse Databricks utilise des tables Delta, tandis que dbt formalise les dépendances, la qualité et le modèle dimensionnel. La table de faits est définie au grain d'une ligne de commande et alimente un schéma en étoile consommable dans Power BI.

L'étude détaille les exigences fonctionnelles et non fonctionnelles, les arbitrages d'architecture, les mécanismes d'idempotence et les tests. Une validation de bout en bout confirme la propagation d'un client, d'une commande et de ses lignes jusqu'aux dimensions et à la table de faits. Le résultat constitue une chaîne reproductible et observable, adaptée à la démonstration de changements quasi continus sans remettre en cause un pipeline batch déjà stable.

\textbf{Mots-clés :} Data Engineering, PostgreSQL, Databricks, Delta Lake, dbt, Apache Airflow, Power BI, traitement incrémental, architecture médaillon, schéma en étoile.

\clearpage
\chapter*{Abstract}
\addcontentsline{toc}{chapter}{Abstract}
This report presents the design and implementation of \emph{Walmart Streaming Flow}, an educational data engineering platform built to simulate, ingest, transform and analyze retail events. A Python and Flask web generator produces two business-consistent event families: automatically generated customer sign-ups and orders containing one to five line items. Writes are committed transactionally to an external PostgreSQL database while preserving foreign-key relationships between customers, stores, employees and products.

The analytical processing intentionally remains an hourly micro-batch workflow. Apache Airflow orchestrates eleven deterministic steps: incremental ingestion into Bronze, freshness checks, dbt transformations into technical and business Silver layers, historical snapshots and construction of the Gold layer. Databricks stores the lakehouse data in Delta tables, while dbt formalizes dependencies, data quality rules and dimensional modeling. The fact table is defined at order-line grain and feeds a Power BI-ready star schema.

This study describes functional and non-functional requirements, architectural trade-offs, idempotency mechanisms and verification strategy. An end-to-end validation confirms that a newly created customer, order and line item are propagated into the appropriate dimensions and fact table. The resulting solution is reproducible and observable, and demonstrates continuously changing source data without destabilizing an already reliable batch transformation pipeline.

\textbf{Keywords:} Data Engineering, PostgreSQL, Databricks, Delta Lake, dbt, Apache Airflow, Power BI, incremental processing, medallion architecture, star schema.
